\section{Introduction}
\label{sec:intro}

North African social-media text is written in Arabic script, in Latin-script Arabizi with digits standing for Arabic consonants (\textit{3}=\ar{ع}, \textit{7}=\ar{ح}), and with French segments inside the same sentence. A classifier trained on Arabic-script data can be applied to Arabizi in two ways: on the original string, or after transliterating it to Arabic script. Transliteration helps when it recovers words the classifier knows, and harms when it garbles French segments, negation, names or the spelling conventions of another dialect. The usual choice is made once for a whole dataset \citep{muller2021unseen,moosa2023transliteration,artetxe2023revisiting}, although the right choice varies from one input to the next.

This paper asks whether a frozen classifier can learn, without labels at inference time, \emph{for which inputs} normalization pays off, and whether that decision transfers from Tunisian to Algerian and Moroccan Arabizi. Learning a per-input benefit predictor is not new: post-hoc deferral rules estimate the gain of a second stage \citep{jitkrittum2023cascade}, and loss predictors choose test-time transformations \citep{kim2020learningloss}. What these methods do not address is how such a rule should be moved to a new dialect, where it is most likely to be wrong and where labels are scarce.

We start from an elementary identity. With binary labels, the expected benefit of normalizing an input equals the probability that the two views disagree, times the margin by which the normalized view wins when they do (\cref{eq:factor}). The first factor requires no labels; the second requires labels only on disagreements. The identity is not new in substance (it underlies disagreement-based methods such as query-by-committee), and we use it as an analysis device: it splits the routing problem into disagreement detection, conditional advantage estimation and a cost-aware decision, which can be estimated, adapted and \emph{measured} separately. We ask whether per-input selection beats dataset-level choices and strong routers (RQ1), across corpora and dialects (RQ2); what it costs (RQ3); whether normalized predictions agree with predictions on human transliterations (RQ4); which sub-problem limits it (RQ5); and how much a few target labels help (RQ6). All conclusions concern frozen encoders with linear heads. We contribute:
\begin{itemize}[leftmargin=1.2em,itemsep=0pt,topsep=2pt]
\item a disagreement-factorized analysis of selective normalization: a cost-aware Bayes rule that separates sunk from avoidable inference cost, and estimators and adaptation regimes for each factor, built from standard tools (\cref{sec:formulation,sec:method});
\item empirical evidence, across four frozen encoders and three Maghrebi varieties, on how far the \emph{conditional advantage} of a normalized view, the factor that per-input routing hinges on, can be predicted from decision-point features with linear and boosted estimators (\cref{sec:mechanism});
\item a controlled, pre-specified evaluation with latency-based costs checked by end-to-end timing, paired tests with multiplicity correction, and interval estimates for the mechanism measurements (\cref{sec:setup,sec:results});
\item a conservative Arabizi normalizer with safeguards for French segments, negation and non-lexical cues; its intrinsic scores are development results, not an independent test (\cref{tab:intrinsic}).
\end{itemize}
Normalization helps substantially on Tunisian and Algerian text, and its sign on Moroccan text depends on the encoder. SelectNorm stays within one point of the strongest router, but neither it nor any learned router consistently beats the better single view, and per-input routing is more expensive than either single view because the first pass is sunk. The analysis locates the limit in the conditional advantage: disagreements are moderately predictable, but the winning view is only weakly predictable from the features and estimators we test, even in domain. We report this rather than tune the method against the test sets.
